\section{A common small gauge away from very large seeds}
\label{sec:seed-tree-patch}
The remaining energy comparison needs a gauge in which both the fine
and coarse fields are small near an unhit plaquette. The fine field
is given on its actual block-tree slice: its tree bonds equal $1$.
We use that tree, including paths crossing a periodic seam.

\subsection{Comparing one fine bond with the coarse replica}
In each block the rooted forest gives a path from a fine site to its
block centre, of length at most $4L$. On the tree slice its transport
is $1$, since every step is a tree bond or its inverse. For a fine
bond $e=(x,a)$, concatenate the reverse path from $x$ to its root,
the bond itself, and the path from $x+e_a$ to its root. This gives a
walk $w$ between the two roots with
\[
 |w|\le8L+1,\qquad U_w=U_e.
\]
Lift this entire walk to the covering lattice starting at zero,
with origin the first root. Its length fits in the macro radius
$(J+7)L$ because $J\ge4$. Its endpoint $z$ lies in that macro box.
For fine period at least $100(J+7)L$, wrapping is injective on the
macro box. Hence $z=0$ if the two roots coincide, and $z=Le_a$ if
the bond crosses to the next block. The latter is true even at a
periodic seam: the periodic endpoints agree with the two stated
cover endpoints, and injectivity identifies their lifts.

Use the macro gauge $H$ of Section~\ref{sec:seed-geometry} at the
first coarse root $i$, and write $P_i$ for its macro plaquette sum.
If its bond sum is at most $D_0LP_i$, the chord of the transport of
$w$ is at most its length times that bond sum. This deliberately
allows repeated steps and gives
\begin{equation}\label{eq:tree-fine-link}
 r(H_0U_eH_z^{-1})\le (8L+1)D_0L P_i.
\end{equation}
The normalized local frames take value $1$ at their own centres.
Thus the centre term of \eqref{eq:geometry-macro-frame} gives
\[
 r(H_{Ly}c_y)\le C_0L^2P_i\quad(y=0,e_a).
\]
Together with the gauged predictor bond bound this yields
\[
 r(H_0Q_{i,a}H_{Le_a}^{-1})
       \le(2C_0L^2+C_1L^{-2})P_i.
\]
Indeed insert $c_0c_0^{-1}$ and $c_{e_a}c_{e_a}^{-1}$ and apply the
triangle inequality to the two centre errors and the predictor term.

The coarse replica $\overline V$ is $1$ on internal fine bonds and
$V_{i,a}$ on crossing bonds. If $e$ is internal,
\eqref{eq:tree-fine-link} has $z=0$ and conjugation removes $H_0$.
If it crosses, compare the last two displayed bounds in the same
macro gauge and then compare $Q_{i,a}$ with $V_{i,a}$. Thus, with
\[
 D_L=(8L+1)D_0L+2C_0L^2+C_1L^{-2},
\]
we obtain the shell comparison
\begin{equation}\label{eq:tree-shell}
 r(U_e\overline V_e^{-1})\le D_LP_i+
 \begin{cases}
 0,& e\text{ internal to a block},\\
 r(V_{i,a}Q_{i,a}^{-1}),&e\text{ crossing from }i.
 \end{cases}
\end{equation}
The constant is allowed to depend on the fixed block side. The only
frame hypothesis is the same positive-weight condition in the macro
comparison. No replacement of the prescribed forest was made.

\subsection{Constructing the gauge in an unhit neighbourhood}
Fix an outer plaquette with anchor $o$ which is not in the hit region.
Let $T$ be the very large threshold. Throughout distance
$\lfloor4B/5\rfloor$ from $o$, all fine plaquette chords and all
predictor discrepancies are at most $T$. For listed defects this
follows from being unhit: any value above $T$ would place a very
large seed in that neighbourhood. For unlisted defects it follows
from the sector bounds $t/2\le T$ and $2Wt\le T$.

Write $M=(J+7)L$ and $K_L=(2M+1)^4N_p$. For $B$ sufficiently large
after fixing $L,W$, every macro box whose coarse centre lies within
$r=\lfloor B/2\rfloor$ of $o$ is contained, in coarse parent
distance, in the above neighbourhood. Its macro sum is at most
$K_LT$, with lifted multiplicities included. The necessary frames
are active there as well. By
\eqref{eq:geometry-predictor-curvature}, every coarse plaquette with
anchor within $r$ has chord at most
\[
 F_LT,\qquad F_L=4+CL^{-2}K_L.
\]
The allowed coarse period is at least $100(1+B)$, so the cover box
of radius $r$ is nonwrapping. Apply the axial gauge construction
\eqref{eq:geometry-axial} on this box. Its plaquette sum is at most
$(2r+1)^4N_pF_LT$, so the total gauged bond sum is at most
$16(2r+1)^5N_pF_LT$. Descend the cover gauge by the centred lift
$y\mapsto\operatorname{lift}(y-o)$. For a bond whose anchor is
within $\lfloor2B/5\rfloor$ of $o$, both endpoints lie in the box;
the descended gauge therefore agrees with the cover gauge at both
endpoints. Each such coarse bond has chord at most
\begin{equation}\label{eq:tree-coarse-small}
 H_L B^5T,\qquad H_L=16\,3^5N_pF_L.
\end{equation}
This follows by bounding one nonnegative summand by the total and
using $2r+1\le3B$.

Call the resulting coarse gauge $z$ and use $z\circ\operatorname{parent}$
on the fine lattice. Under this block-constant gauge, the replica
of $V$ becomes the replica of $z\cdot V$: odd $L$ ensures the
prescribed parent changes across a positive fine bond either not at
all or by the corresponding positive coarse step. Thus its bond
chord is either zero or bounded by \eqref{eq:tree-coarse-small}.
On the other hand, the chord of $U_e\overline V_e^{-1}$ is unchanged
by the simultaneous gauge, since that product is conjugated at the
initial endpoint. Equation \eqref{eq:tree-shell} consequently gives
\[
 r((z\circ\operatorname{parent})\cdot U)_e
       \le(D_LK_L+1)T+H_LB^5T.
\]
With $C_L=H_L+D_LK_L+1$ and $B\ge1$, both fine and coarse fields
therefore satisfy the common bound
\begin{equation}\label{eq:tree-common-gauge}
 r((z\circ\operatorname{parent})\cdot U)_e\le h(x),\qquad
 r(z\cdot V)_b\le h(x),\qquad h(x)=C_LB^{10}T,
\end{equation}
on the radius-$\lfloor2B/5\rfloor$ neighbourhood of $o$.
The fine bound uses parent distance for the anchor. All choices,
including the gauge, refer to the same original fields, tree and
tags. The threshold for $x$ may depend on $L,W$ and the fixed
construction parameters, but not on the torus volume or on the
chosen unhit plaquette. This proves the common small-gauge input
needed for the local comparison.
